\subsection{Long-time decay and N-waves (Evans \S3.4.5)}
For this subsection assume $f(0)=0$, $f$ smooth, and
$0<\theta\leq f''\leq\Theta<\infty$.
These convenient sufficient hypotheses ensure uniform convexity of
both $f$ and $L=f^*$, as required in the source asymptotic arguments.
Write $\sigma=f'(0)$ and $d=f''(0)$.

\begin{theorem}[Supremum decay; source Theorem 5]
For bounded integrable initial data,
$\|u(\cdot,t)\|_\infty\leq C t^{-1/2}$.
\end{theorem}
\begin{expandedproof}
The primitive $P$ in (LOP) is bounded by $M=\|g\|_1$.
Duality gives $L(\sigma)=0$, $L'(\sigma)=0$.
Since $L''\geq1/\Theta$,
$L(v)\geq(v-\sigma)^2/(2\Theta)$.
Compare a minimizing $y$ with $y_0=x-\sigma t$:
\[
 \frac{|x-y-\sigma t|^2}{2\Theta t}-M
 \leq tL((x-y)/t)+P(y)\leq P(x-\sigma t)\leq M.
\]
Thus $|(x-y)/t-\sigma|\leq C/\sqrt t$.
As $G(\sigma)=0$ and $G$ is $1/\theta$-Lipschitz,
formula (LO) gives the claim.
\end{expandedproof}
\Needspace{45mm}
\PDEAnnotation{衰减的是高度，不一定是质量}{波的高度可按 $t^{-1/2}$ 降低，同时宽度按 $\sqrt t$ 增长。积分质量不一定趋零，所以长期形状应和一个展开的 N 波比较，而非只和零比较。}

Suppose additionally $\supp g\subset[-R,R]$.
Use the equivalent primitive $P(x)=\int_{-\infty}^xg$,
so $P_-=0$, $P_+=m=\int g$.
Let $P_*=\min P$,
\[
 p=-2P_*,\qquad q=2(m-P_*),\qquad p,q\geq0.
\]
The associated N-wave is
\[
 N(x,t)=\begin{cases}
 (x-\sigma t)/(dt),
   &-\sqrt{pdt}<x-\sigma t<\sqrt{qdt},\\
 0,&\text{otherwise}.
 \end{cases}
 \tag{NW}
\]
Its negative and positive signed areas are $-p/2$ and $q/2$,
so $\int N=(q-p)/2=m$.
\Needspace{45mm}
\PDEAnnotation{N 波参数由原函数最低点决定}{$p,q$ 不是初值正负质量各乘二。例如先取正脉冲再取等质量负脉冲，原函数始终非负且最终回零，此时 $p=q=0$，即使初值两种符号都有。}
\begin{theorem}[Source Theorem 6]
If $p,q>0$, then $\|u(\cdot,t)-N(\cdot,t)\|_1\leq C t^{-1/2}$.
\end{theorem}
\begin{expandedproof}[locating the edges before integrating the error]
Pass to the moving coordinate $\xi=x-\sigma t$ and put
$\ell(v)=L(\sigma+v)$. Then
$\ell(0)=\ell'(0)=0$ and $\ell''(0)=1/d$.
The minimizing velocity bound just proved is $|(\xi-y)/t|\leq C/\sqrt t$.
Taylor expansion on this fixed small-velocity range gives
\[
 t\ell((\xi-y)/t)=\frac{(\xi-y)^2}{2dt}+O(t^{-1/2}),
\]
\[
 u(x,t)=\frac{\xi-y(x,t)}{dt}+O(t^{-1}).
 \tag{NT}
\]
The constants are uniform for minimizers; the same Taylor bound holds
for the bounded starting-point competitors used below when $|\xi|=O(\sqrt t)$.

Choose $\varepsilon=A/\sqrt t$ with a fixed sufficiently large $A$.
We work for large $t$, so $\varepsilon<1$.

\paragraph{1. Exterior edges: the solution is zero beyond them.}
If $\xi<-R-\sqrt{pd(1+\varepsilon)t}$, the competitor $y=\xi$ has
objective $P_-=0$.
A minimizing $y\geq-R$ would, by (NT), have objective at least
\[
 \frac{pd(1+\varepsilon)t}{2dt}+P_*+O(t^{-1/2})
 =\frac{pA}{2\sqrt t}+O(t^{-1/2})>0.
\]
This contradicts the competitor when $A$ is large enough.
Thus $y<-R$, where $P=0$, and the strictly positive cost $\ell$
is minimized only by $y=\xi$. Hence $u=G(\sigma)=0$.
The right-tail argument is the same with the baseline $P_+=m$ and $q$.
Therefore $u=0$ outside
\[
 -R-\sqrt{pd(1+\varepsilon)t}
 \leq \xi\leq
 R+\sqrt{qd(1+\varepsilon)t}.
 \tag{OUT}
\]

\paragraph{2. Interior edges: the minimizing point stays in the initial support.}
At $\xi_l=R-\sqrt{pd(1-\varepsilon)t}$, any $y<-R$ has objective
at least $P_-=0$. Choose a point $z\in[-R,R]$ with $P(z)=P_*$.
Since $|\xi_l-z|\leq\sqrt{pd(1-\varepsilon)t}$ for large $t$,
its objective is at most
\[
 \frac{p(1-\varepsilon)}2+P_*+O(t^{-1/2})
 =-\frac{pA}{2\sqrt t}+O(t^{-1/2})<0.
\]
Thus every minimizer at $\xi_l$ is at least $-R$.
Similarly at $\xi_r=-R+\sqrt{qd(1-\varepsilon)t}$ it is at most $R$.
Monotonicity of minimizers then gives $|y(x,t)|\leq R$
throughout $\xi_l<\xi<\xi_r$.
There (NT) implies
$|u(x,t)-\xi/(dt)|\leq C/t$.

\paragraph{3. Integrate over the middle and the transition strips.}
The middle interval has length $O(\sqrt t)$ and lies inside the ideal
N-wave interval, so its $L^1$ error is $O(t^{-1/2})$.
Each strip between the interior, ideal, and exterior edges has bounded
length, since
$\sqrt{(1\pm A/\sqrt t)t}-\sqrt t=O(1)$.
Both $|u|$ and $|N|$ are $O(t^{-1/2})$, giving the same error there.
Outside (OUT) both are zero. This proves the bound for large times.

For a bounded time interval the minimizing-velocity bound from (HL)
gives finite propagation of the compact support and $|u|\leq\|g\|_\infty$.
Thus $\|u\|_1$ is uniformly bounded on that interval; $\|N\|_1=(p+q)/2$.
Enlarge $C$ to cover all remaining $t>0$.
\end{expandedproof}
\Needspace{45mm}
\PDEAnnotation{边缘位置也要估计}{只有波内 $O(1/t)$ 点态误差还不够做全空间积分。还要把两侧支集误差压进宽度 $O(1)$ 的条带，再乘以高度 $O(t^{-1/2})$，才能得到总误差。}

For Burgers flow with initial data $\mathbf1_{(0,1)}$, the fan is $u=x/t$.
Until $t=2$ it meets a height-one plateau ending at shock $x=1+t/2$.
After $t=2$ the shock is $x=\sqrt{2t}$ and
$u=x/t$ on $0<x<\sqrt{2t}$, zero elsewhere.
This follows from the fan--shock matching equation
$\gamma'=\gamma/(2t)$ and $\gamma(2)=2$.
It has height $O(t^{-1/2})$ and constant mass one.
Here $p=0$, $q=2$; the explicit example illustrates decay, while
the general $L^1$ theorem above was stated with $p,q>0$.
